\documentclass[11pt]{article}
\usepackage{fontspec}
\usepackage[french,english]{babel}
\usepackage{amsmath}
\usepackage{amssymb}
\usepackage{amsthm}
\usepackage{mathrsfs}
\usepackage{tikz-cd}
\usepackage{enumitem}

\newcommand{\printedpage}[1]{\marginpar{\footnotesize p.~#1}}
\newcommand{\msp}[1]{\marginpar{\footnotesize [#1]}}
\newcommand{\Filt}{\operatorname{Filt}}
\newcommand{\Ob}{\operatorname{Ob}}
\newcommand{\card}{\operatorname{card}}
\newcommand{\Fl}{\operatorname{Fl}}
\newcommand{\Hom}{\operatorname{Hom}}
\newcommand{\uHom}{\underline{\operatorname{Hom}}}
\newcommand{\Ens}{\mathrm{Ens}}
\newcommand{\Ind}{\operatorname{Ind}}
\newcommand{\circled}[1]{\tikz[baseline=(c.base)]\node[draw,circle,inner sep=1pt](c){#1};}

\begin{document}

\setcounter{footnote}{39}

\printedpage{34}
% head: A. Grothendieck : Dérivateurs, Ch. XVIII, Édition des Universités de Montpellier II et Paris VII
% foot: Version 17 décembre 2025

\msp{page 117}

\begin{quote}
\textit{where $\mathcal{M}_\pi$ denotes the strictly full subcategory of $\mathcal{M}$ formed of the $\pi$-accessible objects.}
\end{quote}

\begin{quote}
\textit{This being so, we have the implications}
\[
\mathrm{1)} \ \Longrightarrow\ \mathrm{2)} \ \Longrightarrow\ \mathrm{3)} \ \Longrightarrow\ \mathrm{4)} \ \Longrightarrow\ \mathrm{5)} \ ,
\]
\textit{and in the particular case where $\mathcal{M} = $} (Ens), \textit{$C$ = full subcategory formed of the sets of cardinal $\leqslant \pi_0$, we also have} 5) $\Longrightarrow$ 1), \textit{i.e. all these conditions are equivalent.}
\end{quote}

\textsc{Proof.}

1) $\Longrightarrow$ 2). Cf. proof of 4.5.4.

2) $\Longrightarrow$ 3). The category $J = \Filt^\pi/F$ is stable under $\varinjlim$ of cardinal $\leqslant \pi$, $\Filt^\pi$ being so, hence $J$ is large relative to $\pi$, moreover $F = \varinjlim\limits_{j\in J} X_j$, where $j \longmapsto X_j$ is the source functor $\Filt^\pi/F \longrightarrow \Filt^\pi$. Whence the desired representation of $F$, replacing $J$ if need be by a cofinal ordered set $I$, e.g. by taking for $I$ the set of subcategories $J$ of $I$ such that $\card \Fl J \leqslant \pi$ \texttt{[this second part of the sentence ``e.g. ...'' calls for correction]}. The bound $\card I \leqslant \pi'^{\pi}$, for $F \in \Ob \Filt^{\pi'}$, is obtained by AQT \texttt{[âne qui trotte: at a donkey's trot, i.e. by routine verification]}.

The implication 3) $\Longrightarrow$ 4) is tautological.

4) $\Longrightarrow$ 5) (\footnote{compare lemma 4.5.14, page 133.}), for if $F \in \Ob \mathcal{M}_\pi$, we find by 4) that $F$ is a direct factor of an object of $\Filt^\pi$, and $\Filt^\pi$ being stable under $\varinjlim$
\msp{page 118}
large relative to $\pi_0$, this implies that $\Filt^\pi$ is stable under direct factors, whence $F \in \Ob \Filt^\pi$.

Finally, if $\mathcal{M} = \Ens$, $C = \Ens(\pi_0)$. We note that
\[
\Filt^\pi \ = \ \{ \ X \in \Ob \Ens \mid \card\,(X)^{\pi_0} \ \leqslant \ \pi \ \} \ .
\]
On the other hand, we easily find
\[
\begin{aligned}
(\Ens)_\pi \quad &\overset{\text{def}}{=} \quad \text{full subcategory of (Ens) formed} \\
&\phantom{\overset{\text{def}}{=}} \quad \text{of the $\pi$-accessible objects} \\
&= \quad \{ X \in \Ob \Ens \mid \card X \leqslant \pi \} \ .
\end{aligned}
\]
Thus we have $(\Ens)_\pi \subset \Filt^\pi$ if and only if $\pi^{\pi_0} \leqslant \pi$, i.e. $\pi^{\pi_0} = \pi$, q.e.d. (\footnote{Cf. proof of part 1°) of 4.5.6, below (bottom of page and following page).}).

\textbf{Corollary 4.5.7.} \textit{For every cardinal $\pi \geqslant \pi_1$, we have}
\[
\Filt^\pi \ \subset \ \mathcal{M}_\pi \ \subset \ \Filt^{\pi^{\pi_0}} \ . \tag*{(4.5.7.1)}
\]

The inclusion $\Filt^\pi \subset \mathcal{M}_\pi$ is the implication a $\Longrightarrow$ c of part 1°) of 4.5.6. To prove $\mathcal{M}_\pi \subset \Filt^{\pi^{\pi_0}}$, we put $\pi' = \pi^{\pi_0}$ and note that $\pi'^{\pi_0} = \pi^{\pi_0 \cdot \pi_0} = \pi^{\pi_0} = \pi'$. Thus we can
\printedpage{35}
% head: A. Grothendieck : Dérivateurs, Ch. XVIII, Édition des Universités de Montpellier II et Paris VII
% foot: Version 17 décembre 2025
apply 1° to $\pi'$ instead of $\pi$, and we find that $\mathcal{M}_{\pi'} = \Filt^{\pi'}$. Since $\mathcal{M}_\pi \subset \mathcal{M}_{\pi'}$, we conclude $\mathcal{M}_\pi \subset \Filt^{\pi'}$, q.e.d.

I forgot to prove 1°) in 4.5.6. In view of the diagram of implications $(\ast)$, p. 108, and the implication $\mathrm{a}'' \Longrightarrow \mathrm{a}$ already proved (p. 109),
\msp{page 119}
it remains only to prove
\[
\mathrm{c}' \ \Longrightarrow \ \mathrm{a}'' \ ,
\]
i.e. that $F$ $\pi$-accessible \textit{in} $C^\wedge$ implies $\card F(x) \leqslant \pi \ \forall \ x \in C$. This will follow from the somewhat more general result:

\textbf{Lemma 4.5.8.}
\begin{itemize}[leftmargin=3em]
\item[\textit{1°)}] \textit{In the category $\mathcal{E} = $} (Ens), \textit{the $\pi$-accessible objects ($\pi$ an infinite cardinal) are the sets $X$ such that} $\card X \leqslant \pi$.
\item[\textit{2°)}] (\footnote{Cf. generalization 4.14.18, p. 391 b) and c).}). \textit{Let $C$ be a category equivalent to a small category, such that $\forall \ x, y \in \Ob C$, we have} $\card \Hom(x,y) \leqslant \pi$. \textit{Let $\mathcal{E}$ be a $\mathfrak{U}$-category stable under $\varinjlim$ large relative to $\pi$, and weakly $\pi$-accessible (i.e. admitting a generating subcategory formed of $\pi$-accessible objects), and stable under products of cardinal $\leqslant \pi$. Then for an object $F$ of $\uHom(C,\mathcal{E})$ to be $\pi$-accessible, it is} necessary \textit{that $\forall \ x \in \Ob C$, $F(x)$ be a $\pi$-accessible object of $\mathcal{E}$, and this condition is also sufficient if} $\card \operatorname{\text{red}} \Fl(C) \leqslant \pi$ (\footnote{even \textit{without} any condition of stability of $\mathcal{E}$ under products of cardinal $\leqslant \pi$.}).
\end{itemize}

\textsc{Proof.} 1°) $\card X \leqslant \pi \Longrightarrow X$ $\pi$-accessible in $\mathcal{E} = (\Ens)$, for $X = \coprod_{x \in X} \{x\}$, and we apply SGA 4 I 9.9. Conversely, suppose $X$ $\pi$-accessible, and write $X \simeq \varinjlim\limits_{i \in I} X_i$, where $I$ is the set of subsets of $X$ which are of cardinal $\leqslant \pi$. Since $I$ is large relative to $\pi$, the identity map $X \xrightarrow{\ \mathrm{id}\ } X$ factors through a set $X_i$, hence $X_i = X$, hence $\card X = \card X_i \leqslant \pi$, q.e.d.

\msp{page 120}

2°) Consider, for $x \in \Ob C$, the functor
\[
e \xrightarrow{\ \varepsilon_x\ } C \ ,
\]
whence a functor
\[
\varepsilon_x^* \ : \ \uHom(C,\mathcal{E}) \longrightarrow \mathcal{E} \ , \qquad F \longmapsto F(x) \ ,
\]
admitting a right adjoint
\[
\varepsilon_{x\,*} \ : \ \mathcal{E} \ \longrightarrow \ \uHom(C,\mathcal{E}) \ ,
\]
given by
\[
\varepsilon_{x\,*}(E)(y) \ = \ E^{\Hom(x,y)}
\]
(subject to the existence of products of cardinal $\leqslant \pi$). This functor $\varepsilon_{x\,*}$ is $\pi$-accessible by virtue of loc. cit. 9.8.

This being so, let $F$ be a $\pi$-accessible object of $\uHom(C,\mathcal{E})$, we see that $\varepsilon_x^*(F)$ is $\pi$-accessible, $\varepsilon_x^*$ having a $\pi$-accessible right adjoint:
\[
E \ \longmapsto \ \Hom_{\mathcal{E}}(\underbrace{F(x)}_{\varepsilon_x^*(F)}, E) \ \simeq \ \Hom_{\uHom(C,\mathcal{E})}(F, \varepsilon_{x\,*}(E))
\]
\printedpage{36}
% head: A. Grothendieck : Dérivateurs, Ch. XVIII, Édition des Universités de Montpellier II et Paris VII
% foot: Version 17 décembre 2025
is the composite of the $\pi$-accessible functor $\varepsilon_{x\,*}$, and the $\pi$-accessible covariant functor represented by $F$ in $\uHom(C,\mathcal{E})$.

Conversely, if the $F(x)$ are $\pi$-accessible, let us prove that $F$ is so, provided $\card \Fl(C) \leqslant \pi$. Done by AQT \texttt{[âne qui trotte: at a donkey's trot, i.e. by routine verification]}, I am not going to spell it out here.

\msp{page 121}

\textbf{Proposition 4.5.9.} \textit{Let $C$ be a category equivalent to a small category, and let $\mathcal{M} = $} $\Ind'_{\pi_0}(C)$, \textit{$\pi_0$ a given infinite cardinal. Assume $C$ stable under $\varinjlim$ of cardinal $\leqslant \pi_0$. For every cardinal $\pi \geqslant \pi_0$, let $\mathcal{M}_\pi \supset C$ be the full subcategory of $\mathcal{M}$ formed of the $\pi$-accessible objects of $\mathcal{M}$. Let, on the other hand,} $\pi_1 = \sup(\pi_0, \card \Fl C)$. \textit{Then $(\mathcal{M}_\pi)_{\pi \geqslant \pi_1}$ is a cardinal filtration of $\mathcal{M}$} (\footnote{\textbf{N.B.} All the conditions for cardinal filtrations are satisfied with the exception, when $\pi \geqslant \pi_0$, of that of the bound on cardinals $\card I \leqslant \pi'^{\pi}$ condition c).})\textit{, and moreover, for $\pi \geqslant \pi_0$, $\mathcal{M}_\pi$ is stable in $\mathcal{M}$ under $\varinjlim$ of cardinal $\leqslant \pi$.}

\textsc{Proof.} The stability of $\mathcal{M}_\pi$ is a reminder, cf. SGA 4 I 9.9. Since we have
\[
\mathcal{M}_\pi \ \subset \ \Filt^{\pi^{\pi_0}}
\]
for $\pi \geqslant \pi_1$ (cf. 4.5.7), we therefore see that the $\mathcal{M}_\pi$ are equivalent to small categories (condition a) for cardinal filtrations). Condition b) for cardinal filtrations is contained in the stability property of the $\mathcal{M}_\pi$ which we have just recalled. Finally let us prove condition c). So let $\pi \geqslant \pi_0$, and let $x$ be in $\mathcal{M}$, consider the category
\[
I \ = \ \mathcal{M}_\pi/x \ ,
\]
it is stable under $\varinjlim$ of cardinal $\leqslant \pi$, by the similar stability of $\mathcal{M}_\pi$, hence $I$ is large relative to $\pi$ and equivalent to a small category, on the other hand, we have
\msp{page 122}
\[
x \ = \ \varinjlim_{i \in I} x_i \ ,
\]
the $x_i$ in $\mathcal{M}_\pi$ ( $\varinjlim$ in $\mathcal{M}$ of the functors $I = \mathcal{M}_\pi/x \longrightarrow \mathcal{M}_\pi$). Replacing $I$ if need be by a cofinal ordered category, this proves the qualitative part of c. Suppose finally $x \in \Ob \mathcal{M}_{\pi'}$, where $\pi' > \pi$, let us show that we then have
\[
\card I_0, \ \card I' \ \leqslant \ \pi'^{\pi} \qquad \text{(\footnotemark, \footnotemark)} \ , \tag*{(4.5.9.1)}
\]
\addtocounter{footnote}{-1}%
\footnotetext{where $I_0$, $I'_0$ are skeletal subcategories of $I$, $I'$ such that $I_0 \hookrightarrow I$, $I'_0 \hookrightarrow I'$ are equivalences \texttt{[Grothendieck has crossed out the subscript 0 of $I'$ in 4.5.9.1]}.}%
\stepcounter{footnote}%
\footnotetext{We only prove it if $\pi \geqslant \pi_1$.}%
which will finish proving that we have a cardinal filtration. We need a bound on the cardinals.

\textbf{Lemma 4.5.10} (\footnote{Give a common statement with 4.5.11, and first prove 4.5.11.}). \textit{Let $\pi$ be a cardinal $\geqslant \pi_1$. Then we have} % [sic] énonce -> énoncé, missing accent (in footnote 47)
\[
\card \Fl((\Filt^\pi)_0) \ \leqslant \ 2^\pi \ . \tag*{(4.5.10.1)}
\]

Indeed we have, denoting by $C^\wedge(\pi)$ the strictly full subcategory of $C^\wedge$ formed of the $F : C^\circ \longrightarrow \Ens$ with $\card F(x) \leqslant \pi$ for every $x$ in $C$ (i.e. $F$ factors through $\Ens(\pi)$), and by $C^\wedge(\pi)_0$ an
\printedpage{37}
% head: A. Grothendieck : Dérivateurs, Ch. XVIII, Édition des Universités de Montpellier II et Paris VII
% foot: Version 17 décembre 2025
\msp{page 123}
exhaustive skeletal subcategory,
\[
\card C^\wedge(\pi)_0 \ \leqslant \ 2^\pi \ , \tag*{(4.5.10.2)}
\]
whence (4.5.10.1), since for $\pi \geqslant \pi'$, $\Filt^\pi$ is a full subcategory of $C^\wedge(\pi')$. To establish (4.5.10.2), we note that the set of isomorphism classes of $\Ens(\pi')$ is identified with the set of cardinals $\leqslant \pi'$, which set is of cardinal $\leqslant \pi'$. Replacing $\Ens(\pi)$ by $\Ens(\pi')_0$, $C$ by $C_0$, we must bound $\card \Fl(\uHom(C_0, \Ens(\pi')_0))$. Now $\Ob \uHom(C_0, \Ens(\pi')_0)$ is contained in the set of maps from $\Fl C_0$ to $\Fl(\Ens(\pi')_0)$, the first is bounded by $\pi_1$, the second by $\pi^\pi \times \pi \times \pi = \pi^{2\pi} = \pi^\pi = 2^\pi$ (\textbf{N.B.} $\pi < 2^\pi$, whence $\pi^\pi \leqslant (2^\pi)^\pi = 2^{\pi^2} = 2^\pi$), hence $\Ob \uHom$ bounded by $(2^\pi)^{\pi_1} = 2^{\pi\pi_1} = 2^\pi$. On the other hand, if $F$, $G$ are in $\uHom(C_0, \Ens(\pi_0))$, $\Hom(F,G)$ is bounded by $(\pi^\pi)^{\pi_1} = \pi^\pi = 2^\pi$, hence $\Fl \uHom$ is bounded by $2^\pi \times 2^\pi \times 2^\pi$, q.e.d.

\textbf{Corollary 4.5.11.} \textit{Let $F \in \Ob \Filt^\pi$, $F' \in \Ob \Filt^{\pi'}$, with $\pi' \geqslant \pi \geqslant \pi_1$. Then}
\[
\card \Hom(F, F') \ \leqslant \ \pi'^{\pi} \ . \tag*{(4.5.11.1)}
\]

\msp{page 124}

It is essentially the same proof, slightly simpler. (We should have started there.)

\textbf{Corollary 4.5.12.} \textit{If $\pi' \geqslant \pi \geqslant \pi_1$, we have for $x \in \Ob \mathcal{M}_{\pi'}$}
\[
\card\,(\mathcal{M}_\pi/x)_0 \ \leqslant \ \pi'^{(\pi^{\pi_0})} \ . \tag*{(4.5.12)}
\]

Indeed,
\[
\card \Ob\,(\mathcal{M}_\pi/x)_0 \ \overset{\text{by 4.5.11}}{\leqslant} \ \underbrace{\underbrace{\card \Ob\,(\mathcal{M}_\pi)_0}_{\leqslant 2^{(\pi^{\pi_0})} \text{by 4.5.10}} \times (\pi'^{\pi_0})^{\pi^{\pi_0}}}_{= \pi'^{(\pi^{\pi_0})}} \ ,
\]
i.e.
\[
\card \Ob\,(\mathcal{M}_\pi/x)_0 \ \leqslant \ \pi'^{(\pi^{\pi_0})} \ ,
\]
whence
\[
\card \Fl(\mathcal{M}_\pi/x)_0 \ \leqslant \ \pi'^{(\pi^{\pi_0})} \times \pi'^{(\pi^{\pi_0})} \times (\pi'^{\pi_0})^{(\pi^{\pi_0})} \ \leqslant \ \pi'^{(\pi^{\pi_0})} \ .
\]
This therefore gives (returning to 4.5.9.1)
\[
\card I_0 \ \leqslant \ \pi'^{(\pi^{\pi_0})} \ .
\]

On the other hand, $I'$ can be defined in terms of $I_0$ as the set of subcategories of $I_0$ which are of cardinal $\leqslant \pi$, hence their number is bounded by $(\card I_0)^\pi \ \leqslant \ (\pi'^{\pi^{\pi_0}})^\pi = \pi'^{\pi\pi^{\pi_0}} = \pi'^{(\pi^{\pi_0})}$.

But this is not quite the required bound, since
\[
\pi'^{\pi} \leqslant \pi'^{(\pi^{\pi_0})} \qquad \text{(\footnotemark)},
\]
\footnotetext{\textbf{N.B.} We have equality if we have $\pi' = 2^{c'}$ with $c' \geqslant \pi^{\pi_0}$, or $\pi = 2^c$ with $c \geqslant \pi_0$.}%
and since it is possible, a priori, that the inequality
\printedpage{38}
% head: A. Grothendieck : Dérivateurs, Ch. XVIII, Édition des Universités de Montpellier II et Paris VII
% foot: Version 17 décembre 2025
\msp{page 125}
is strict! Apparently, the ``bounds''
\[
\mathcal{M}_\pi \ \subset \ \Filt^{\pi^{\pi_0}} \ , \qquad \mathcal{M}_{\pi'} \ \subset \ \Filt^{\pi'^{\pi_0}}
\]
are too crude. So it is a matter of getting a better grasp of the categories $\mathcal{M}_\pi$, for $\pi \geqslant \pi_0$ (\footnote{But we can use 4.5.10, 4.5.11 to conclude that for $x \in \Ob \mathcal{M}_{\pi'} \subset \Ob \Filt^{\pi'^{\pi_0}}$, we have $\card \Fl(C_0/x) \leqslant (\pi'^{\pi_0})^{\pi_0} = \pi'^{\pi_0}$, we will therefore have $x = \varinjlim\limits_{J} x_j$ ($x_j \in \Ob C$, $\varinjlim$ filtering, $\card \Fl J \leqslant \pi'^{\pi_0}$), and introducing the set $I$ of subcategories $J' \subset J$ which are of cardinal $\leqslant \pi$, and putting $x_{J'} = \varinjlim\limits_{j \in J'} x_j$, we find $x = \varinjlim\limits_{J' \in I} x_{J'}$, where this time $I$ is an \textit{ordered set} large relative to $\pi$, the $x_{J'} \in \mathcal{M}_\pi$, and $\card I \leqslant (\card J)^\pi \leqslant (\pi'^{\pi_0})^\pi = \pi'^{\pi\pi_0} = \pi'^{\pi}$, this is the right bound!}). I am going to venture a

\textbf{Proposition 4.5.13.} \textit{The hypotheses and notations are those of 4.5.9. Let $\pi$ be a cardinal $\geqslant \pi_0$, and let $x \in \Ob \mathcal{M}$. Consider the conditions:}
\begin{itemize}[leftmargin=3em]
\item[\textit{a})] \textit{There exists a category $I$, with} $\card \Fl I \leqslant \pi$, \textit{and a functor $i \longmapsto x_i : I \longrightarrow C$, such that $x$ is isomorphic to $\varinjlim x_i$.}
\item[\textit{a}$'$)] \textit{As in a), with $I$ filtering.}
\item[\textit{a}$''$)] \textit{As in a), with $I$ ordered and filtering.}
\item[\textit{b})] \textit{There exists a} cofinal \textit{functor $I \longrightarrow C/x$, with} $\card \Fl I \leqslant \pi$.
\item[\textit{b}$'$)] \textit{The filtering category $C/x$ is ``essentially of cardinal $\leqslant \pi$'', i.e. there exists an ordered category} \texttt{[no, cf. b$''$)]} filtering $I$, $\card \Fl I \leqslant \pi$, \textit{and a cofinal functor $I \longrightarrow C/x$.}
\item[\textit{b}$''$)] \textit{As in b$'$), with $I$ moreover ordered.}
\item[\textit{c})] \textit{There exists a strict epimorphic family of arrows $(u_\alpha)_{\alpha \in A}$, $x_\alpha \longrightarrow x$, indexed by a set $A$ such that} $\card A \leqslant \pi$.
\item[\textit{d})] \textit{We have $x \in \Ob \mathcal{M}_\pi$, i.e. $x$ is $\pi$-accessible in $\mathcal{M}$.}
\end{itemize}

\msp{page 126}

\textit{This being so:}
\begin{itemize}[leftmargin=3em]
\item[\circled{1\textsuperscript{\textit{o}}}] \textit{The first six conditions $a$, $a'$, $a''$, $b$, $b'$, $b''$ are equivalent, say $P(\pi)$, and imply the last two,}
\[
\begin{tikzcd}[column sep=small]
 & P(\pi) \arrow[dl, Rightarrow] \arrow[dr, Rightarrow] & \\
c & & d \ .
\end{tikzcd}
\]
\textit{If $\pi \geqslant \pi_1$, then $P(\pi)$ is equivalent to $c$.}
\item[\circled{2\textsuperscript{\textit{o}}}] \textit{Let $\Filt'^{\pi} \subset \mathcal{M}$ be the strictly full subcategory of $\mathcal{M}$ formed of the objects which satisfy condition $P(\pi)$ above, and $\Filt^\pi$ the strictly full subcategory defined in 4.5.3 (p. 109). We then have}
\[
\Filt^\pi \ \subset \ \Filt'^{\pi} \ \subset \ \mathcal{M}_\pi \ , \tag*{(4.5.13.1)}
\]
\textit{and for there to be equality of the three terms, it suffices that we have $\pi^{\pi_0} = \pi$. (Cf. prop. 4.5.6, page 116). Moreover, $\Filt'^{\pi}$ is stable in $\mathcal{M}$ under} sums \textit{of cardinal $\leqslant \pi$.}
\item[\circled{3\textsuperscript{\textit{o}}}] \textit{The following conditions on $C$, $\pi_0$, $\pi$ are equivalent.}
\end{itemize}

\end{document}
